\documentclass{hhuarticle}
% \usepackage[ngerman,english]{babel} % English
\usepackage[english,ngerman]{babel} % Deutsch

\title{Gestalten der schriftlichen Ausarbeitung}

% Vorlange von Jannik Dunkelau und Joshua Schmid
\author{John Witulski} % keine Matrikelnummer eintragen

\subject{Bachelorseminar: Programmiersprachen} % Optional.
\semester{Wintersemester 2022}

% \printblackwhitetrue % Vereinfacht Schwarzweißdruck (Logo, Plots, ...).

%%%%%%%%%%%%%%%%%%%%%%%%%%%%%%%%%%%%%%%%%%%%%%%%%%%%%%%%%%%%%%%%%%%%%%%%%%%%%%%%
%% LaTeX Packages in Nutzung                                                  %%
%%                                                                            %%
%% Im folgenden können Sie für die Niederschrift benötigte LaTeX-Pakete       %%
%% einbinden.                                                                 %%
%% Diese Vorlage kommt bereits mit einigen nützlichen inkludierten Paketen.   %%
%%%%%%%%%%%%%%%%%%%%%%%%%%%%%%%%%%%%%%%%%%%%%%%%%%%%%%%%%%%%%%%%%%%%%%%%%%%%%%%%

%% Häufig benutzte mathematische Packages.
\usepackage{amsfonts}
\usepackage{amsmath}
\usepackage{amssymb}

\usepackage{enumitem} % Leichter konfigurierbare enumerate-Umgebungen.
\usepackage{subcaption} % Unterteilung von Figures in Subfigures.
\usepackage[colorlinks]{hyperref} % Klickbare Links (z.B. Inhaltsverzeichnis).
\sethyperrefhhucolors{} % Setzt den Farbsatz der HHU für hyperref.
\usepackage{url} % \url Kommando für Darstellung von Links
\usepackage{csquotes} % Improved quoting.

%% Tabellen
\usepackage{tabularx} % tabularx Umgebung für mehr Kontrolle über Tabellen.
\usepackage{booktabs} % \toprule, \midrule, \bottomrule
\usepackage{multirow}
\usepackage{multicol}
\usepackage{longtable} % Große Tabellen gehen über mehrere Seiten.

%% Quellcode
\usepackage{listings} % Einbindung von Code.
\setlstlistingstyle{} % Kosmetische Einstellungen
% Sprachabhängige Bezeichnung.
\iflanguage{ngerman}{\renewcommand{\lstlistingname}{Quellcode}}{}

%% Algorithmen in Pseudocode
\usepackage{algorithm} % Float-Umgebung für angegebene Algorithmen.
\usepackage{algorithmicx} % Angabe von Algorithmen in Pseudocode.
\usepackage{algpseudocode} % Standart Pseudocode-Elemente für Algorithmen.
\algsmallfont{}
% Sprachabhängige Bezeichnung.
\iflanguage{ngerman}{\floatname{algorithm}{Algorithmus}}{}

%% Intelligenteres Referenzieren mittels \cref.
\usepackage[capitalize,noabbrev,ngerman]{cleveref}
%\pagenumbering{gobble}

\definecolor{pblue}{rgb}{0.13,0.13,1}
\definecolor{pgreen}{rgb}{0,0.5,0}
\definecolor{pred}{rgb}{0.9,0,0}
\definecolor{pgrey}{rgb}{0.46,0.45,0.48}
\lstset{language=Java,
  showspaces=false,
  showtabs=false,
  breaklines=true,
  showstringspaces=false,
  breakatwhitespace=true,
  commentstyle=\color{pgreen},
  keywordstyle=\color{pblue},
  stringstyle=\color{pred},
  basicstyle=\ttfamily,
  moredelim=[il][\textcolor{pgrey}]{$$},
  moredelim=[is][\textcolor{pgrey}]{\%\%}{\%\%}
}

\begin{document}

  \maketitle

  \begin{abstract}
    % Was ist das Thema
    Haskell ist eine Sprache, welche vor mehr als 30 Jahren zuerst erschien
    und dennoch selbst heutzutage ihrer Zeit in manchen Bereichen voraus ist.
    % Warum ist das Thema interessant
    Trotz dessen ist sie in der Industrie wenig verbreitet, obwohl
    Programmierer viel von ihr lernen können.
    % Wie wurde es umgesetzt
    Inwiefern Programmierer davon profitieren können,
    wird in dieser Arbeit untersucht, indem
    Parallelen zwischen Haskell und anderen Sprachen gezogen werden.
    % Welche Ergebnisse wurden erzielt
    Wir erkennen dabei, dass Haskell mächtige Strukturen besitzt, welche
    in anderen Sprachen nicht oder nur teilweise existieren, da Haskell
    Datenstrukturen auf einem viel abstrakteren Level behandelt,
    wodurch Probleme mit Haskell weitaus effizienter und eleganter
    gelöst werden können.
    % Welche Gründe dies hat und warum sich dies ändern sollte, insbesondere
    % mit Blick auf Programmierer 
    % nicht nur Haskell, sondern
    % auch andere Sprachen, sowie Programmierer davon profitieren können,
    % wird in dieser Arbeit erläutert.
  \end{abstract}

\tableofcontents


  \section{Einleitung}

  Im September 1987 wurde die
  ''Conference on Functional Programming Languages and Computer Architecture``
  in Portland, Oregon, USA, abgehalten. Ziel war es, die zersplitterte
  funktionale Programmiersprachen-Community zu vereinen und eine
  gemeinsame Standardsprache zu entwickeln.
  Die Sprache sollte alle Vorteile der bis dahin existierenden
  funktionalen Sprachen, insbesondere der zu der Zeit populären Sprache
  ''Miranda``, vereinen und gleichzeitig die Nachteile beseitigen.
  Die Sprache, welche dabei entstand, ist Haskell.

  Während ihrer Entwicklung wurde Haskell stark von Akademikern beeinflusst,
  wobei insbesondere Themen wie Typentheorie, Kategorientheorie und das Lambda-Kalkül
  eine große Rolle spielten. Daraus resultierte, dass Haskell
  einige Besonderheiten aufweist, welche es von anderen Sprachen unterscheidet.
  Diese Besonderheiten sind es, welche Haskell zu einer so interessanten
  Sprache machen, da somit Haskell eine der wenigen reinen
  Programmiersprachen ist, also keine Seiteneffekte vorhanden sind.
  Eine unintuitive Konsequenz daraus ist, dass Haskell keine Variablen besitzt,
  welche verändert werden können, da es keine Zuweisung gibt,
  da Zuweisungen in Haskell Seiteneffekte erzeugen würden. Somit
  existieren in Haskell nur Konstanten, welche nicht verändert werden können.
  Dies wiederum führt dazu, dass klassische Schleifen in Haskell nicht existieren,
  da diese durch Seiteneffekte, insbesondere dem Inkrementieren einer
  Laufvariable, implementiert werden. Stattdessen werden in Haskell
  Funktionen verwendet, welche rekursiv aufgerufen werden, um
  die gewünschte Schleifenfunktionalität zu erreichen.

  Statt der Objektorientierung, verwendet Haskell die funktionale
  Programmierung, was durch die Reinheit bedeutet, dass Haskell
  durch allein kategorientheoretische Konzepte definiert ist.
  ``Kategorientheoretische Konzepte`` beziehen sich auf die
  Kategorientheorie, welche eine Theorie ist, welche
  eine alternative Grundlage für die Mathematik darstellt, im Gegensatz
  zu der klassischen Mengenlehre.
  Diese besagt, dass es Kategorien gibt, in welchen es sogenannte Objekt
  gibt, welche durch Morphismen verbunden sind. Existieren
  Objekte $A$, $B$ und $C$ in einer Kategorie $\mathcal{C}$,
  sowie Morphismen $f: A \to B$ und $g: B \to C$, so gibt es einen
  Morphismus $g \circ f: A \to C$ und außerdem muss es für jedes
  Objekt $A$ einen Morphismus $id_A: A \to A$ geben, welcher
  das Objekt $A$ nicht verändert. In dem Fall von Haskell sind die
  Objekte die Typen und die Morphismen die Funktionen, wobei dies noch
  weitreichende Konsequenzen hat, welche im Folgenden noch aufgefasst werden.
  Funktionale Programmierung bedeutet, dass Programme durch Funktionen
  definiert werden, welche Daten als Eingabe erhalten und Daten als
  Ausgabe liefern. Geschickte Komposition von Funktionen ermöglicht
  es, komplexe Probleme durch hintereinander aufrufende einfache
  Funktionen zu lösen.
  Dies hat zur Folge, dass Programmierer, welche Haskell lernen,
  automatisch besser in abstrakter Mathematik werden. Außer der Reinheit
  und der funktionalen Programmierung, ist Haskell eine
  polymorphe, statisch-typisierte Programmiersprache mit
  lazy-evaluation, wobei diese Begriffe im Folgenden erläutert werden.

  Polymorphie bedeutet, dass Funktionen und Datenstrukturen
  generisch sind, was bedeutet, dass sie mit verschiedenen Typen verwendet werden können.
  Statische Typisierung bedeutet, dass der Typ einer Variable
  zur Kompilierzeit festgelegt wird, wodurch Fehler, welche durch
  dynamische Typisierung entstehen, vermieden werden. Durch das
  algebraische Typensystem von Haskell, was bedeutet, dass Typen $A$ und
  $B$ zu Produkttypen $A \times B$ und Summentypen $A | B$ kombiniert werden können,
  ist es möglich, Funktionen und Datenstrukturen mit verschiedenen
  Typen zu kombinieren. Hier sehen wir wieder, dass ein starker Fokus
  auf dem Kombinieren von einfacheren Strukturen liegt, wodurch komplexere
  Strukturen entstehen. Da Haskell auch eine funktionale Programmiersprache ist,
  haben auch Funktionen einen Typen, wodurch Funktionen, Konstanten und Variablen
  alle durch die starke Typisierung von Haskell definiert werden können
  und somit eine Typenstruktur aufweisen. Da wir bereits wissen, dass
  Haskell eine stark-typisierte Programmiersprache ist, bedeutet dies,
  dass Programmierfehler in Haskell fast immer von dem Compiler erkannt werden.
  Sollte eine Funktion mit einem falschen Typen aufgerufen werden, oder
  einen falschen Typen zurückgeben, so wird dies vom Compiler erkannt.
  Diese Eigenschaft von Haskell ist sehr nützlich, da es so möglich ist,
  sicherheitskritische Programme zu schreiben, welche mit erhöhter
  Wahrscheinlichkeit korrekt sind als Programme in anderen Sprachen.

  Die lazy-evaluation von Haskell bedeutet, dass Ausdrücke erst dann
  ausgewertet werden, wenn sie benötigt werden. Dies hat zur Folge,
  dass Ausdrücke, welche nie benötigt werden, nicht ausgewertet werden.
  Dies ist sehr nützlich, da es so möglich ist, unendliche Datenstrukturen
  zu definieren, welche nur dann ausgewertet werden, wenn sie benötigt werden.
  Beispielsweise kann eine Liste von allen Primzahlen definiert werden,
  wobei man nun die $n$-te Primzahl aus dieser Liste erhalten kann,
  ohne dass die Liste vollständig ausgewertet werden muss.
  Auch wird Code somit idiomatischer, da man analog zur natürlichen
  Sprache ''die $n$-te Primzahl`` verwenden kann, statt diese
  durch eine Funktion zu erhalten, welche die $n$-te Primzahl zurückgibt.

  Diese Eigenschaften machen Haskell zu einer sehr mächtigen Sprache, welche
  es ermöglicht, elegante und effiziente Lösungen für viele Probleme zu finden.
  Durch ihren starken mathematisch, theoretischen Grundaufbau besitzt Haskell
  Strukturen tief in sich verankert, welche noch heutzutage
  von anderen Programmiersprachen teilweise übernommen werden.
  Wie genau diese Strukturen in Haskell implementiert sind, wird im Folgenden
  erläutert. Außerdem werden wir sehen, wie Spezialisierungen 


  % Jede Ausarbeitung beginnt mit einer guten Einleitung.
  % Hier liegt der Fokus auf den Grundlagen,
  % sodass klar ist, worum es geht, was das Thema behandelt und
  % wieso dies in erster Linie interessant ist.
  % Detaillierte Informationen folgen in den späteren Abschnitten.

  % Diese grundlegende Einleitung erlaubt auch,
  % dass spätere Informationen hierauf aufbauen können.
  % Dadurch lässt sich verhindern,
  % dass zwei unterschiedliche Abschnitte die gleichen Grundlagen definieren
  % und somit im Dokument doppeln.


  \section{Struktureller Aufbau}

  Es ist ein zusammenhängender Text von 15--20 Seiten (+ Referenzen) zu verfassen,
  der das gegebene Thema und die Lernziele abdeckt.
  Achten Sie hierbei auf eine sinnvolle Gliederung der Unterthemen.

  \subsection{Unterthemen}
  Unterthemen geben bereits eine Aufteilung in Abschnitte vor und jedes
  Unterthema sollte nicht mehr als einen Abschnitt umfassen.
  Dabei sollte nach dem Titel des Abschnitts je ein
  kurzer, einführender Text stehen, bevor der Abschnitt in weitere,
  sinnvolle Unterabschnitte aufgeteilt wird.
  Dieser einführende Text steht direkt hinter dem jeweiligen Titel und
  nicht in einem etwaigen Unterabschnitt.


  \section{Abbildungen, Tabellen und Beispiele}

  Sie werden in Ihren Zusammenfassungen nicht darum herum kommen,
  Abbildungen, Tabellen und Beispiele zur besseren Anschaulichkeit zu nutzen.
  Diese Abbildungen müssen im Text referenziert und besprochen werden
  und sollten nicht alleine stehen.
  Ebenfalls sollten Abbildungen und Tabellen zentriert dargestellt werden
  (mit \texttt{\textbackslash centering}).

  Mit \cref{fig:initial-draft} wird beispielhaft die Referenzierung einer Grafik
  demonstriert.
  Auf \LaTeX-Seite geschieht dies durch das Setzen eines Labels
  (\texttt{\textbackslash label}) für die Grafik
  und durch die Nutzung der Befehle
  \texttt{\textbackslash ref} oder \texttt{\textbackslash cref}.

  \begin{figure}[h]
    \centering
    \includegraphics[width=4cm]{fig/hhulogo.pdf}
    \caption{Das neue HHU-Logo.}%
    \label{fig:initial-draft}
  \end{figure}

  Der \texttt{\textbackslash cref}-Befehl entstammt hierbei dem
  \texttt{cleveref}-Paket und setzt automatisch den korrekten Bezeichner
  vor die Referenz, in diesem Beispiel ``Abbildung''.
  Zum Vergleich: \texttt{\textbackslash ref} setzt nur die
  Nummer \ref{fig:initial-draft}, die der Abbildung zugeteilt wurde.

  \subsection{Beispiele}
  Viele Themen profitieren von einem guten Beispiel,
  um Konzepte oder Eigenschaften klarer zu präsentieren.
  Achten Sie bei Beispielen darauf,
  dass diese nicht unnötig komplex sind und nur auf das nötige reduziert sind.
  Sollte es sich anbieten, ist es hilfreich ein Running Example zu nutzen,
  welches von allen Unterthemen aufgegriffen wird.
  So ist der Leser bereits mit den Gegebenheiten des Beispiels vertraut
  und kann Unterschiede zwischen den verschiedenen Szenarien genauer
  differenzieren.


  \subsection{Code}

  \begin{figure}[h]

\begin{lstlisting}
public class HelloWorld 
{
 
       public static void main (String[] args)
       {
             // Ausgabe Hello World!
             System.out.println("Hello World!");
       }
}
\end{lstlisting}
 \caption{So binden Sie Code ein}%
  \label{fig:codeBlock}
\end{figure}

  \subsection{Tabellen}

\begin{table}[h!]
\centering
\begin{tabular}{ |c|c|c| } 
 \hline
 Sprache & Jahr & Version \\ 
\hline
 Java & 1997 & JSE 19.0 \\ 
 C & 1972 & C17 \\ 
 \hline
\end{tabular}
\caption{Eine sinnfreie Tabelle}
\label{table:1}
\end{table}

  \section{Literaturverweise}

  Die Angabe von Quellen und Verweisliteratur ist essentiell in Ihrer
  Ausarbeitung,
  als auch später in Ihrer Bachelorarbeit.

  Literaturverweise werden in \LaTeX{} mittels
  \texttt{\textbackslash cite} getätigt.
  Verweise stehen hinter der Aussage, die sie stützen sollen,
  oder hinter dem Namen eines referenzierten Konzepts oder Eigenschaft
  und gehören in den Satz hinein.
  Beachten Sie folgendes Beispiel:
  \begin{itemize}
    \item In ``Artificial Intelligence: A Modern Approach''~\cite{russell2002artificial}
      werden moderne Ansätze und Techniken
      der künstlichen Intelligenz vorgestellt.
    \item bei direkten Zitationen wir zusätzlich die Seite angegeben wo das Zitat zu finden ist, z.B. ``...''~\cite[p.1]{russell2002artificial}
  \end{itemize}

  Am Ende der Ausarbeitung werden die referenzierten Quellen automatisch
  aufgelistet.

  \subsection{Bibtex-Einträge und \texttt{references.bib}}

  Um mit \texttt{\textbackslash cite} referenzieren zu können,
  müssen die entsprechenden Einträge in der Datei \texttt{references.bib}
  stehen.
  Sie finden in der Vorlage bereits den beispielhaften Eintrag für
  ``Artificial Intelligence: A Modern Approach''~\cite{russell2002artificial}.

  Für Literaturquellen empfiehlt es sich in aller Regel
  \url{https://scholar.google.com} zu nutzen.
  Hier kann man sich direkt den Bibtex-Eintrag einer Quelle
  angeben lassen und muss diesen so nicht selbst zusammenstellen.

  Beachten Sie auch den Artikel
  ``Common Errors in Bibliographies''
  von John Owens.\footnote{%
    \url{https://www.ece.ucdavis.edu/~jowens/biberrors.html}}
  Des Weiteren ist der Vortrag von Simon Peyton Jones \href{https://www.microsoft.com/en-us/research/academic-program/write-great-research-paper/}{``How to write a great research paper''} zu empfehlen.

  \appendix
  \section{Beispielhafter Anhang}
  Der Anhang wird mit dem \texttt{\textbackslash{}appendix}-Befehl eingeleitet.
  Alle folgenden Sektionen werden als Anhang formatiert.

\backmatter
\listoffigures
\listoftables
\bibliography{references}
\bibliographystyle{plaindin}

\end{document}

% Quellen
%
% Why You’ll Probably Never Use Haskell in Production
% https://henrydangprg.com/2019/05/27/why-youll-probably-never-use-haskell-in-production/
%
% Choosing Haskell isn’t a stand-in for good software design
% https://ozataman.medium.com/choosing-haskell-isnt-a-stand-in-for-good-software-design-7d893882f963
%
% A History of Haskell: Being Lazy With Class
% https://dl.acm.org/doi/10.1145/1238844.1238856
